% 来源：问题分析_Word版.docx。
\section{问题分析}
\subsection{问题一的分析}
问题一是后续各问的基础，其核心是描述预热平衡阶段药材内部的热量与水分传递过程。仅考虑径向变化能够描述径向中心与表面的差异，但无法反映端面交换及轴向梯度。因此，将药材视为二维轴对称区域，同时保留径向和轴向的传热传质。预热阶段药材的密度、比热容、导热系数按题目附录2取定值，水分扩散系数随局部含水率变化；烘房温度和水分浓度由附件1确定。

该问的求解重点是在变化的环境条件下计算温度场和含水率场，并获得指定时刻、位置的结果。本文采用有限体积法进行空间离散，利用自适应BDF方法进行时间积分，由此建立后续问题共用的计算框架。

计算得到整个二维区域内的场量后，选取轴向中截面$z=0$进行展示。温度场和含水率场关于该截面对称，该截面距离两个端面最远，受端面换热和水分交换的直接影响相对较弱，能够反映药材内部的变化。因此，四个问题均提取$T(r,0,t)$和$C(r,0,t)$作为代表性输出。

\subsection{问题二的分析}
在问题一的基础上，采用题目附录3给出的物性关系，将研究范围扩展到初始$3\,\mathrm{h}$烘干阶段。药材密度、比热容和导热系数随局部含水率变化，水分扩散系数同时受到温度和含水率的影响。含水率下降改变传热参数，温度变化又影响扩散速度，使两个物理场相互关联。因此，在积分过程中根据当前状态更新相应系数，同步求解温度场和含水率场。

\subsection{问题三的分析}
问题三要求确定药材达到烘干标准所需的总时间，其传热传质机理与问题二相同，但需要将计算持续推进至药材内部各处均满足烘干要求。为此，持续监测全域最大含水率，并以其首次降至阈值的时刻作为总烘干时间。

附件1仅给出了前$4\,\mathrm{h}$的烘房环境数据，而题目指出实际烘干通常持续2至3天，长期求解必须明确处理数据截止后的环境条件。本文将末20组成对观测的温度和水分浓度均值作为后续环境条件，并通过改变平均窗口检查烘干时间对该选择的敏感程度。

\subsection{问题四的分析}
问题四进一步考虑药材尺寸变化，需要在更新本问密度、比热容、导热系数和水分扩散系数的同时，根据附件2插值得到半径函数$R(t)$。假设药材长度保持不变，仅在径向均匀收缩，将当前径向位置按半径归一化，使节点始终对应药材内部的同一相对位置。计算中只需更新节点实际位置、控制体体积和界面面积，无需重新生成网格连接关系。

烘干终点继续采用全域最大含水率首次达到规定阈值的判据。最后，将本问的烘干时间与问题三比较，分析径向收缩和物性变化对烘干过程的综合影响。
